\documentclass[11pt]{article}

\usepackage[margin=1in]{geometry}
\usepackage{amsmath,amssymb}
\usepackage{graphicx}
\usepackage{booktabs}
\usepackage{hyperref}
\usepackage[numbers]{natbib}

\title{Critical error rate analysis for MSA-based variant phasing\\in nanopore amplicon sequencing}
\author{Josh Walker}
\date{}

\begin{document}
\maketitle

\begin{abstract}
Amplicon consensus tools that detect intra-specimen sequence variants must
distinguish real biological variation from sequencing error. We present a
statistical framework for evaluating the significance of variants detected by
MSA-based position phasing in Oxford Nanopore amplicon data. Rather than
computing $p$-values against a fixed error rate, we derive the \emph{critical
error rate} $p^*$: the per-position sequencing error rate at which an observed
variant pattern would become plausible under a null hypothesis of independent
errors. Using conservative assumptions---total input reads as the sample
denominator and Bonferroni correction over all amplicon positions---we show that
variants with moderate read support require per-position error rates exceeding
those of current nanopore chemistry to be explained as artifacts. We present
results under both a conservative error model ($q = p$, maximally biased
substitution) and a uniform model ($q = p/3$), and examine the effect of
significance level on variant confidence across a range of operating conditions.
We additionally provide minimum read-support thresholds at assumed error rates,
enabling direct lookup of the evidence required for variant confidence.
\end{abstract}

\section{Introduction}

Oxford Nanopore sequencing of PCR amplicons is widely used for taxonomic
identification, particularly in fungal DNA barcoding of the internal transcribed
spacer (ITS) region~\cite{schoch2012}. After demultiplexing, reads from a single
specimen are clustered and consensus sequences are generated to produce
high-accuracy reference sequences suitable for database matching.

When a specimen contains multiple sequence variants---arising from heterozygosity,
paralogous loci, mixed templates, or multi-copy rDNA polymorphism~\cite{bradshaw2023}---the clustering
and consensus pipeline must resolve these into separate consensus sequences rather
than collapsing them into a single ambiguous output. Speconsense~\cite{speconsense} addresses this
through MSA-based variant phasing: after initial clustering, it examines the
partial order alignment~\cite{lee2003,vaser2017} within each cluster to identify
positions where a substantial fraction of reads consistently differ from the
consensus, then recursively splits the cluster into haplotype groups based on
allele patterns at these positions.

A natural concern is whether detected variants represent real biological
sequences or artifacts of sequencing error. If nanopore errors at a particular
position happened to affect the same base in enough reads, the phasing algorithm
could mistake this for a genuine variant. We address this concern by deriving the
per-position error rate required to make each variant plausible as an artifact,
and showing that this critical rate exceeds realistic nanopore error for variants
with sufficient read support.

\section{Methods}

\subsection{Null model}

Under the null hypothesis, all reads in a cluster originate from a single true
sequence, and observed base differences are independent sequencing errors. At
each position, each read independently has probability $p$ of containing an
error.

We present results under two substitution models:

\begin{itemize}
\item \textbf{Conservative} ($q = p$): All errors at a position produce the same
alternative base. This represents the worst case, accounting for nanopore
substitution biases that favor certain error types~\cite{delahaye2021,wick2023}. If errors are
concentrated toward one alternative base, the probability of multiple reads
independently producing the same incorrect base is maximized.

\item \textbf{Uniform} ($q = p/3$): Errors are spread equally across the three
alternative bases. This is more realistic for positions without strong substitution
bias, and provides a less conservative bound.
\end{itemize}

\subsection{Parameter definitions}

\begin{itemize}
\item $N$: Total input read count for the specimen, before clustering. This is
conservative because it represents the full pool from which a biased subsample
could be drawn. In practice, speconsense uniformly presamples to $N \leq 1000$
reads, so this is bounded.

\item $M$: Number of reads supporting the variant allele at a position.

\item $L$: Amplicon length in base pairs. For the ITS region, $L \approx 700$~bp.
This determines the multiple testing correction.

\item $K$: Number of variant positions distinguishing a given variant from all
other variants produced from the same specimen. We use $K = 1$ (single-position
variants) throughout as the conservative bound (Section~\ref{sec:multiposition}).

\item $\alpha$: Significance level. We examine $\alpha = 0.05$ (per-variant),
$\alpha = 10^{-3}$, and $\alpha = 10^{-5}$ (appropriate when correcting for
thousands of variant calls across a sequencing run).
\end{itemize}

\subsection{Critical error rate for single-position variants ($K = 1$)}

For a single position where $M$ of $N$ reads share the same alternative base,
the probability under the null model is given by the binomial survival function:
\begin{equation}
P(\text{position is artifact}) = P\bigl(\text{Binom}(N,\, q) \geq M\bigr)
\end{equation}
where $q = p$ (conservative) or $q = p/3$ (uniform).

With Bonferroni correction over $L$ independent positions, the probability that
\emph{any} position exhibits a variant-like pattern is bounded by:
\begin{equation}
P(\text{any artifact variant}) \leq L \cdot P\bigl(\text{Binom}(N,\, q) \geq M\bigr)
\end{equation}

The critical error rate $p^*$ is defined as the value of $p$ satisfying:
\begin{equation}
\label{eq:critical}
L \cdot P\bigl(\text{Binom}(N,\, q^*) \geq M\bigr) = \alpha
\end{equation}
where $q^* = p^*$ under the conservative model or $q^* = p^*/3$ under the
uniform model. This is a single-variable equation solved numerically via root-finding on the
log-scale survival function. If $p^*$ substantially exceeds the actual
per-position error rate of the sequencing platform, the variant cannot plausibly
be explained by error.

\subsection{Choice of significance level}

A single variant call at $\alpha = 0.05$ has a 5\% false positive probability
under the null. In a typical sequencing run producing thousands of variant calls
across hundreds of specimens, the run-level false positive rate can be
substantial. For a run of $S$ specimens producing an average of $V$ variants
each, the expected number of false positives is $\alpha \cdot S \cdot V$.

With $S = 1000$ and $V \approx 3$, $\alpha = 0.05$ yields ${\sim}150$ expected
false positives per run. A Bonferroni correction to the run level suggests
$\alpha \approx 0.05 / 3000 \approx 1.7 \times 10^{-5}$. We therefore present
results at $\alpha = 10^{-5}$ as a practical run-corrected threshold alongside
the per-variant $\alpha = 0.05$.

\subsection{Multi-position variants ($K > 1$)}
\label{sec:multiposition}

When a variant is distinguished from other variants at $K > 1$ positions, the
null hypothesis requires that the \emph{same} $M$ reads independently produce the
matching error at \emph{all} $K$ positions. Under the assumption of independent
errors across positions, the probability that a single read matches the variant
allele at all $K$ positions is $q^K$, where $q$ is the per-position error
probability. The critical error rate equation becomes:
\begin{equation}
\label{eq:critical_multiK}
\binom{L}{K} \cdot P\bigl(\mathrm{Binom}(N,\, (q^*)^K) \geq M\bigr) = \alpha
\end{equation}
where $(q^*)^K$ is the joint per-read probability of matching the variant at all
$K$ positions, and the Bonferroni correction is over $\binom{L}{K}$ possible
position combinations. No correction over alternative bases (e.g., $3^K$) is
needed because the alternative alleles are observed, not hypothesized: we
condition on the specific bases seen in the variant reads.

The effect is dramatic. At $N = 1000$, $M = 100$, $\alpha = 0.05$ (conservative
model): $K = 1$ gives $p^* = 6.75\%$, while $K = 2$ gives $p^* = 24.19\%$---a
$3.6{\times}$ increase. The combinatorial correction grows from $L = 700$ to
$\binom{700}{2} = 244{,}650$ (a $350{\times}$ increase), but the joint error
probability $q^2$ shrinks so rapidly that it overwhelmingly dominates. Under
production settings ($p = 1.5\%$ assumed error, uniform model, $\alpha = 10^{-5}$,
$N = 1000$), the minimum $M$ for significance drops from 23 at $K = 1$ to 6 at
$K = 2$ and 4 at $K = 3$.

The cross-position independence assumption is strongest when the variant positions
are distant in the amplicon sequence. Positions within the same k-mer context may
exhibit correlated error profiles~\cite{delahaye2021}, weakening the $q^K$ bound.
This is not a reason to avoid the $K > 1$ analysis, but a caveat: for variants
whose distinguishing positions span different sequence regions, the joint model is
well-supported; for positions clustered within a few bases, the effective exponent
may be less than $K$.

The analysis that follows uses $K = 1$ throughout as the conservative bound.
Multi-position variants are substantially more significant than this bound
suggests.

\section{Results}

\subsection{Critical error rates at fixed $N$}

Figure~\ref{fig:vs_M} shows the critical error rate $p^*$ as a function of
variant read count $M$ at $N = 1000$ (the standard presample size) for $K = 1$.
Curves are shown for three significance levels under the conservative model
($q = p$), with the uniform model ($q = p/3$) overlaid at $\alpha = 0.05$ and
$\alpha = 10^{-5}$ for comparison.

Under the conservative model at $\alpha = 0.05$, the critical error rate crosses
the 5\% reference line at approximately $M = 80$, and $M = 100$ yields
$p^* = 6.8\%$. Even at the stringent run-corrected $\alpha = 10^{-5}$, $M = 100$
gives $p^* = 5.5\%$---still well above realistic nanopore error rates. The
difference between significance levels is modest at high $M$, reflecting the
exponential sensitivity of the binomial test.

Under the uniform model, the same cases yield $p^* = 20.3\%$ and $16.6\%$
respectively, providing comfortable margins.

At low $M$, the test's power drops rapidly. The boundary where variants become
questionable depends on $N$, but for typical coverage ($N = 100$--$1000$),
variants with fewer than ${\sim}20$--$30$ supporting reads under the conservative
model yield $p^*$ values in the range of realistic nanopore error rates (1--5\%),
making them difficult to distinguish from artifacts.

\begin{figure}[htbp]
\centering
\includegraphics[width=\textwidth]{figures/fig1_critical_rate_vs_M.pdf}
\caption{Critical error rate $p^*$ vs.\ variant read count $M$ at $N = 1000$,
$K = 1$, $L = 700$. Solid lines: conservative model ($q = p$) at three
significance levels. Dashed lines: uniform model ($q = p/3$) at $\alpha = 0.05$
and $\alpha = 10^{-5}$.}
\label{fig:vs_M}
\end{figure}

\subsection{Effect of total read count}

Figure~\ref{fig:by_N} shows $p^*$ as a function of $M$ for several values of $N$
under both error models ($\alpha = 0.05$). For a given $M$, larger $N$ yields
lower $p^*$ because the same absolute count represents a smaller fraction of the
total pool.

The side-by-side comparison illustrates the roughly threefold difference between
error models. Under the uniform model, even relatively modest support ($M = 20$
at $N = 1000$) yields $p^*$ above 1\%. Under the conservative model, $M \approx
30$ is needed to reach the same threshold.

\begin{figure}[htbp]
\centering
\includegraphics[width=\textwidth]{figures/fig2_critical_rate_by_N.pdf}
\caption{Critical error rate $p^*$ vs.\ variant read count $M$ for several total
read counts $N$ ($K = 1$, $L = 700$, $\alpha = 0.05$). Left: conservative model
($q = p$). Right: uniform model ($q = p/3$).}
\label{fig:by_N}
\end{figure}

\subsection{Summary of key values}

Table~\ref{tab:key_values} summarizes critical error rates for representative
parameter combinations under both error models.

\begin{table}[htbp]
\centering
\caption{Critical error rates $p^*$ for representative parameter combinations
($K = 1$, $L = 700$). Values shown for both conservative ($q = p$) and uniform
($q = p/3$) error models.}
\label{tab:key_values}
\begin{tabular}{rrrrrr}
\toprule
 & & & \multicolumn{2}{c}{$p^*$} & \\
\cmidrule(lr){4-5}
$N$ & $M$ & $\alpha$ & $q = p$ & $q = p/3$ & Note \\
\midrule
1000 & 100 & 0.05     & 6.75\%  & 20.25\% & Standard operating point \\
1000 & 100 & $10^{-5}$ & 5.53\%  & 16.59\% & Run-corrected \\
100  & 10  & 0.05     & 2.18\%  & 6.55\%  & Moderate coverage \\
100  & 10  & $10^{-5}$ & 0.83\%  & 2.50\%  & Run-corrected, moderate coverage \\
50   & 5   & 0.05     & 0.86\%  & 2.57\%  & Boundary case \\
50   & 5   & $10^{-5}$ & 0.15\%  & 0.44\%  & Run-corrected boundary \\
\bottomrule
\end{tabular}
\end{table}

\subsection{Minimum read support at assumed error rates}

While Table~\ref{tab:key_values} shows the error rate required to explain a given
variant as artifact, practitioners often need the converse: given an assumed
platform error rate, how many reads must support a variant for it to be
considered trustworthy? Table~\ref{tab:min_M} inverts the analysis to answer this
directly.

\begin{table}[htbp]
\centering
\caption{Minimum variant read count $M$ required for significance at assumed
per-position error rates ($K = 1$, $L = 700$). Upper panel: conservative model
($q = p$). Lower panel: uniform model ($q = p/3$).}
\label{tab:min_M}
\small
\begin{tabular}{r rrr rrr rrr}
\toprule
 & \multicolumn{3}{c}{$p = 1\%$}
 & \multicolumn{3}{c}{$p = 3\%$}
 & \multicolumn{3}{c}{$p = 5\%$} \\
\cmidrule(lr){2-4} \cmidrule(lr){5-7} \cmidrule(lr){8-10}
$N$ & 0.05 & $10^{-3}$ & $10^{-5}$
    & 0.05 & $10^{-3}$ & $10^{-5}$
    & 0.05 & $10^{-3}$ & $10^{-5}$ \\
\midrule
\multicolumn{10}{l}{\emph{Conservative ($q = p$)}} \\[2pt]
10   &  4 &  5 &  6$^\dagger$ &  5 &  6$^\dagger$ &  7$^\dagger$ &  5 &  7$^\dagger$ &  8$^\dagger$ \\
20   &  4 &  5 &  7 &  6 &  7 &  9 &  7 &  9 & 10 \\
50   &  6 &  7 &  9 &  9 & 11 & 13 & 11 & 13 & 15 \\
100  &  7 &  9 & 11 & 12 & 15 & 17 & 16 & 19 & 22 \\
200  & 10 & 12 & 15 & 18 & 21 & 24 & 25 & 28 & 32 \\
500  & 16 & 19 & 22 & 32 & 37 & 41 & 46 & 52 & 57 \\
1000 & 25 & 29 & 33 & 53 & 59 & 65 & 79 & 86 & 93 \\
\midrule
\multicolumn{10}{l}{\emph{Uniform ($q = p/3$)}} \\[2pt]
10   &  3 &  4 &  5 &  4 &  5 &  6$^\dagger$ &  4 &  5 &  6$^\dagger$ \\
20   &  3 &  4 &  5 &  4 &  5 &  7 &  5 &  6 &  8 \\
50   &  4 &  5 &  7 &  6 &  7 &  9 &  7 &  9 & 10 \\
100  &  5 &  6 &  8 &  7 &  9 & 11 &  9 & 11 & 13 \\
200  &  6 &  8 & 10 & 10 & 12 & 15 & 13 & 16 & 18 \\
500  &  9 & 11 & 14 & 16 & 19 & 22 & 22 & 26 & 29 \\
1000 & 13 & 16 & 19 & 25 & 29 & 33 & 35 & 40 & 44 \\
\bottomrule
\end{tabular}

\smallskip
\noindent$^\dagger$\small{$M > N/2$: the required read support exceeds half the
total reads, making detection of multiple variants impossible at this sample
size.}
\end{table}

At the standard operating point ($N = 1000$) under the conservative model, a
variant needs $M \geq 93$ reads to be significant at the run-corrected
$\alpha = 10^{-5}$ when the assumed error rate is 5\%, but only $M \geq 33$ at
1\% assumed error. Under the uniform model, these thresholds drop to 44 and 19
respectively. At moderate coverage ($N = 100$), even the most stringent
combination (5\% error, $\alpha = 10^{-5}$, conservative) requires only
$M \geq 22$, representing 22\% of total reads.

\subsection{Effect of multi-position evidence on minimum read support}

Table~\ref{tab:min_M_K} shows how the minimum $M$ for significance changes when
the variant is distinguished at $K > 1$ positions. Under the joint model
(Section~\ref{sec:multiposition}), requiring the same reads to error at all $K$
positions dramatically reduces the evidence needed.

\begin{table}[htbp]
\centering
\caption{Minimum variant read count $M$ at $K = 1$, $2$, and $3$ variant
positions (uniform model $q = p/3$, $\alpha = 10^{-5}$, $p = 1.5\%$, $L = 700$).}
\label{tab:min_M_K}
\begin{tabular}{rrrr}
\toprule
$N$ & $K = 1$ & $K = 2$ & $K = 3$ \\
\midrule
10   &  5 &  3 &  3 \\
20   &  6 &  3 &  3 \\
50   &  7 &  4 &  3 \\
100  &  9 &  4 &  3 \\
200  & 11 &  4 &  3 \\
500  & 16 &  5 &  3 \\
1000 & 23 &  6 &  4 \\
\bottomrule
\end{tabular}
\end{table}

At $K = 2$, the minimum $M$ drops to 3--6 reads across the full range of $N$; at
$K = 3$, it is 3--4 reads regardless of sample size. These thresholds are at or
below the minimum count filters typically applied in practice (e.g.,
\texttt{--min-size~5}), meaning that for multi-position variants, the statistical
significance gate is effectively redundant---any variant that survives basic
count filtering is already significant. This underscores how conservative the
$K = 1$ analysis is: variants distinguished at two or more positions need
essentially no additional statistical protection.

\section{Discussion}

\subsection{Interpreting the critical error rate}

The critical error rate $p^*$ represents the per-position error rate at which a
given variant pattern becomes statistically plausible as an artifact. When the
conservative model is used ($q = p$), $p^*$ can be interpreted as the rate of
maximally biased error at a single position---the worst case where all errors
produce the same incorrect base. When the uniform model is used ($q = p/3$),
$p^*$ represents the overall per-position error rate under the assumption that
errors are unbiased across alternative bases.

For Oxford Nanopore reads, per-position substitution error rates typically
range from 1--5\% for simplex reads (Q13--Q20) and below 1\% for duplex or
consensus-corrected reads. Modern R10.4.1 chemistry achieves raw read accuracy of
98.7--99.4\% (Q20+), corresponding to per-position error rates of
0.6--1.3\%~\cite{wang2023}, making the margins above even more comfortable for
well-supported variants. Variants with $p^*$ substantially above these ranges
are robust; variants with $p^*$ near or below these ranges warrant caution.

\subsection{Conservative framing}

Several deliberate conservatisms ensure that critical error rates represent lower
bounds on variant significance:

\begin{enumerate}
\item \textbf{Denominator choice.} Using the total input read count $N$ rather
than the cluster size assumes the worst case: the clustering algorithm could have
drawn any subset from the full pool. In practice, coarse clustering first
separates clearly unrelated sequences (e.g., different loci or contaminant
reads), reducing the effective pool. However, finer-grained cluster boundaries
are themselves determined by the variant signal under evaluation, so using
cluster size as the denominator would be circular. The analysis uses total input
reads to avoid this dependency.

\item \textbf{Maximally biased substitution.} The conservative model ($q = p$)
assumes all errors at a position produce the same alternative base. If errors
were spread across multiple alternatives ($q = p/3$), coherent error patterns
would be less likely.

\item \textbf{Bonferroni correction.} Treating all $L$ positions as independent
tests is conservative; nearby positions in amplicon sequences are not truly
independent, so the effective number of tests is smaller than $L$.

\item \textbf{Minimum $K$.} Using $K = 1$ represents the weakest case; most
variants are distinguished at multiple positions. Because the same reads must
error at all $K$ positions independently, significance increases dramatically
with $K$ (Section~\ref{sec:multiposition}).
\end{enumerate}

\subsection{Homopolymer normalization}

The variant detection algorithm includes homopolymer normalization, which
subtracts homopolymer length variation from the base composition at each position
before evaluating variant thresholds. This removes the largest source of
systematic, position-correlated nanopore error from consideration---approximately
47\% of nanopore errors are homopolymer-related, with deletions dominating over
mismatches and insertions~\cite{delahaye2021}. The null model
presented here does not account for this additional filtering, making the analysis
further conservative: the effective error rate after homopolymer normalization is
lower than the raw per-position rate.

\subsection{Relationship to existing variant calling frameworks}

Prior work on low-frequency variant detection has used binomial and
Poisson-binomial models to test whether an observed allele count is consistent
with sequencing error at an assumed rate. LoFreq~\cite{wilm2012} uses per-read
Phred quality scores to construct a position-specific Poisson-binomial null
distribution, applying Bonferroni correction across genomic positions.
SNVer~\cite{wei2011} uses an aggregate binomial model with Simes' procedure for
multiplicity control. Our framework inverts the question: rather than testing
significance at a fixed error rate, we derive the critical error rate $p^*$ at
which a variant becomes plausible as artifact. This inversion is both more
conservative---it does not depend on the accuracy of per-read quality score
calibration, which can be systematically overestimated for nanopore
data~\cite{delahaye2021}---and more interpretable for practitioners, who can
directly compare $p^*$ against known platform error rates.

\subsection{Environmental and mixed specimens}

The analysis assumes that $N$ reads originate from sequencing a single organism,
as is typical in specimen-based amplicon barcoding. For environmental DNA samples
or specimens containing multiple organisms, the interpretation of $N$ requires
modification. When the read pool contains sequences from multiple distinct
organisms, the relevant $N$ for evaluating intra-organism variants is the number
of reads from the target organism, not the total. The upstream clustering step
partially addresses this by separating reads from different organisms into
distinct clusters, but the formal guarantee applies within the assumption that
the initial cluster contains reads from a single source.

\subsection{Limitations}

The null model assumes independence of errors across reads at each position.
Nanopore errors exhibit well-documented sequence-context dependencies---certain
k-mer contexts systematically produce higher error rates and biased substitution
patterns~\cite{delahaye2021}. The homopolymer normalization step addresses the
largest class of such correlation, and the conservative model ($q = p$) accounts
for substitution directionality by assuming all errors at a position produce the
same alternative base. However, neither addresses the possibility that
context-dependent effects elevate the per-position error rate at specific sites
substantially above the platform average. The critical error rate framework is
designed to accommodate this uncertainty: rather than assuming a fixed error
rate, it reports $p^*$ and leaves the reader to judge whether that rate is
plausible given their platform, chemistry, and sequence context. A variant with
$p^* = 6.8\%$ is robust if the worst-case positional error rate is 2\%, but not
if context effects at that position could plausibly reach 7\%. As nanopore
chemistry continues to improve, both mean error rates and context-dependent
profiles change, making fixed thresholds unreliable; $p^*$ remains interpretable
against whatever platform characterization is current.

This framework addresses bioinformatic artifacts---sequences that never
physically existed in the sequenced library. PCR errors introduced early in
amplification produce sequences that physically exist in the library, and are
therefore out of scope by design: passing the critical error rate test
establishes physical presence in the library, not biological reality. For
taxonomic identification, the primary mitigation against PCR artifacts and other
non-biological sequences is robust reference databases where variants have been
replicated across independent specimens.

\section{Conclusion}

The critical error rate framework provides a transparent basis for evaluating
confidence in detected variants. Under typical operating conditions ($N = 1000$,
$M = 100$), critical rates of 5.5--6.8\% (conservative) or 16.6--20.3\%
(uniform) are well above realistic nanopore error, even at run-corrected
significance levels. At lower coverage ($N = 100$, $M = 10$), the conservative
model yields $p^* = 0.8$--$2.2\%$, approaching the boundary of
distinguishability for lower-quality reads. The framework distinguishes physical
library presence from bioinformatic hallucination; biological significance of
detected variants requires downstream validation against reference databases.
Table~\ref{tab:min_M} provides direct lookup of the minimum read support needed
at assumed error rates of 1--5\%, enabling practitioners to assess variant
confidence without computation.

\subsection*{Acknowledgments}

Substantial portions of the text were drafted with the assistance of Claude
(Anthropic, Opus~4.6). The core ideas, analytical framing, and mathematical
approach were provided by the author; the AI contributed to exposition, \LaTeX{}
formatting, and iterative drafting under human direction. The author takes full
responsibility for the content.

\begin{thebibliography}{99}

\bibitem[Bradshaw et~al., 2023]{bradshaw2023}
M.~J. Bradshaw, U.~Braun, B.~Castlebury, and D.~L. Hawksworth.
\newblock Extensive intragenomic variation in the internal transcribed spacer
region of fungi.
\newblock \emph{iScience}, 26(8):107317, 2023.

\bibitem[Delahaye and Nicolas, 2021]{delahaye2021}
C.~Delahaye and J.~Nicolas.
\newblock Sequencing {DNA} with nanopores: Troubles and biases.
\newblock \emph{PLoS ONE}, 16(10):e0257521, 2021.

\bibitem[Lee, 2003]{lee2003}
C.~Lee.
\newblock Generating consensus sequences from partial order multiple sequence
alignment graphs.
\newblock \emph{Bioinformatics}, 19(8):999--1008, 2003.

\bibitem[Schoch et~al., 2012]{schoch2012}
C.~L. Schoch, K.~A. Seifert, S.~Huhndorf, V.~Robert, J.~L. Spouge, C.~A.
Levesque, W.~Chen, and Fungal Barcoding Consortium.
\newblock Nuclear ribosomal internal transcribed spacer ({ITS}) region as a
universal {DNA} barcode marker for fungi.
\newblock \emph{Proceedings of the National Academy of Sciences},
109(16):6241--6246, 2012.

\bibitem[Walker, 2024]{speconsense}
J.~Walker.
\newblock Speconsense: High-quality clustering and consensus generation for
{Oxford Nanopore} amplicon reads.
\newblock \url{https://github.com/joshuaowalker/speconsense}, 2024.

\bibitem[Vaser et~al., 2017]{vaser2017}
R.~Vaser, I.~Sovi\'{c}, N.~Nagarajan, and M.~\v{S}iki\'{c}.
\newblock Fast and accurate de novo genome assembly from long uncorrected reads.
\newblock \emph{Genome Research}, 27(5):737--746, 2017.

\bibitem[Wang et~al., 2023]{wang2023}
Y.~Wang, Y.~Zhao, A.~Bollas, Y.~Wang, and K.~F. Au.
\newblock The newest {Oxford Nanopore} {R10.4.1} full-length {16S} {rRNA}
sequencing enables the accurate resolution of species-level microbial community
profiling.
\newblock \emph{Applied and Environmental Microbiology}, 89(10):e0060523, 2023.

\bibitem[Wei et~al., 2011]{wei2011}
Z.~Wei, W.~Wang, P.~Hu, G.~J. Lyon, and H.~Hakonarson.
\newblock {SNVer}: a statistical tool for variant calling in analysis of pooled
or individual next-generation sequencing data.
\newblock \emph{Nucleic Acids Research}, 39(19):e132, 2011.

\bibitem[Wick et~al., 2023]{wick2023}
R.~R. Wick, L.~M. Judd, and K.~E. Holt.
\newblock Assembling the perfect bacterial genome using {Oxford Nanopore} and
{Illumina} sequencing.
\newblock \emph{PLOS Computational Biology}, 19(3):e1010905, 2023.

\bibitem[Wilm et~al., 2012]{wilm2012}
A.~Wilm, P.~P.~K. Aw, D.~Bertrand, G.~H.~T. Yeo, T.~Ong, C.~H. Wong,
C.~C. Khor, R.~Petric, M.~L. Hibberd, and N.~Nagarajan.
\newblock {LoFreq}: a sequence-quality aware, ultra-sensitive variant caller for
uncovering cell-population heterogeneity from high-throughput sequencing
datasets.
\newblock \emph{Nucleic Acids Research}, 40(22):11189--11201, 2012.

\end{thebibliography}

\end{document}
